\documentclass[11pt,a4paper]{article}

\usepackage[margin=2.5cm]{geometry}
\usepackage{amsmath,amssymb}
\usepackage{bm}
\usepackage{graphicx}
\usepackage{hyperref}

\title{Minutes of Discussion:\\
Thermal Dense Associative Memory, Noise, and Retrieval Phase Transitions}
\author{}
\date{\today}

\begin{document}

\maketitle

\section*{Context of the Discussion}

The discussion focused on clarifying the conceptual role of noise and temperature in associative memory models, and on positioning our current work relative to classical spin-based Hopfield models and recent continuous DAM formulations.

The starting point was the standard applied interpretation of associative memory: stored patterns are corrupted by noise at inference time, and the system dynamics is expected to recover the original pattern. Temperature and noise are often introduced to justify robustness, convergence, or retrieval complexity.

Our goal was to contrast this perspective with the thermodynamic and geometric viewpoint adopted in our work.

%==============================================================================
\section{Classical Spin-Based Interpretation}
%==============================================================================

In classical Hopfield-type models (including polynomial Hopfield networks):

\begin{itemize}
    \item Stored patterns are discrete spin configurations ($\pm 1$).
    \item A corrupted input is generated by flipping a fraction $\gamma$ of spins.
    \item Retrieval is defined as a dynamical process that restores the original pattern.
    \item Temperature and noise are used as physically motivated parameters affecting stability and convergence.
\end{itemize}

This formulation provides a clear applied justification: the system explicitly models recovery from corrupted inputs.

However, it remains fundamentally tied to a \emph{dynamical} and \emph{algorithmic} notion of retrieval.

%==============================================================================
\section{Our Setting and Modeling Choices}
%==============================================================================

We confirmed that our analysis starts from a similar Hamiltonian formulation, but differs in several essential aspects.

\subsection{Continuous State Space}

\begin{itemize}
    \item The system state is continuous and constrained to the sphere of radius $\sqrt{N}$.
    \item Discrete spins are replaced by continuous variables.
    \item Modern Dense Associative Memory kernels are used (LSE, LSR, etc.).
\end{itemize}

These kernels naturally fit into a thermodynamic formalism and admit analytic expressions for internal and free energy.

\subsection{Primary Object of Analysis}

Rather than simulating retrieval dynamics from noisy inputs, we focus on:

\begin{itemize}
    \item analytical derivation of the free energy,
    \item temperature as a control parameter,
    \item structural changes in the free-energy landscape.
\end{itemize}

The emphasis is therefore geometric and thermodynamic, not algorithmic.

%==============================================================================
\section{Macroscopic Description via Overlaps}
%==============================================================================

To analyze the free energy, microscopic states are reparameterized through overlaps with stored patterns:
\[
\phi_\mu = \langle \bm{x}, \bm{\xi}^\mu \rangle.
\]

Here:
\begin{itemize}
    \item $\phi_\mu \approx 1$ corresponds to alignment with pattern $\mu$,
    \item $\phi_\mu \approx 0$ corresponds to states unrelated to stored memories.
\end{itemize}

The free energy is expressed in terms of these macroscopic variables. For interpretability, low-dimensional sections (one or two overlaps) are analyzed.

%==============================================================================
\section{Free Energy Landscape and Temperature}
%==============================================================================

The discussion highlighted the following qualitative behavior:

\begin{itemize}
    \item At low temperature, local minima of the free energy are centered near stored patterns.
    \item As temperature increases:
    \begin{itemize}
        \item minima shift toward the $\phi \approx 0$ region,
        \item basins of attraction shrink,
        \item distinct minima may merge.
    \end{itemize}
\end{itemize}

This behavior is interpreted as a degradation of memory structure rather than failure of a retrieval algorithm.

Temperature controls the \emph{geometry} of memory, not merely stochasticity in the input.

%==============================================================================
\section{Conceptual Role of Temperature}
%==============================================================================

A key point of discussion was the distinction between:

\begin{itemize}
    \item noise in the input patterns,
    \item temperature appearing in the free energy.
\end{itemize}

The latter is interpreted as an intrinsic property of the memory substrate:
\begin{itemize}
    \item internal fluctuations,
    \item imperfect isolation,
    \item analog or physical realization of memory.
\end{itemize}

Thus, temperature is not equivalent to input noise. It limits the very existence and separability of memory states.

%==============================================================================
\section{Phase Transitions and Memory Degradation}
%==============================================================================

The natural scientific objective emerging from the discussion is the study of phase transitions:

\begin{itemize}
    \item control parameters: temperature and memory load,
    \item observables: disappearance or merging of free-energy minima,
    \item interpretation: transition from functional to non-functional memory.
\end{itemize}

Retrieval failure is therefore characterized as a geometric and thermodynamic phase transition.

%==============================================================================
\section{High-Capacity DAMs and Geometric Interference}
%==============================================================================

An important observation concerns modern DAMs:

\begin{itemize}
    \item LSE/LSR kernels support exponential memory capacity,
    \item overlaps $\{\phi_\mu\}$ cannot remain independent in fixed-dimensional space,
    \item strong geometric interference between memories is unavoidable.
\end{itemize}

This further motivates a geometric analysis rather than classical dynamical recovery arguments.

%==============================================================================
\section{Positioning and Scope}
%==============================================================================

Based on the discussion:

\begin{itemize}
    \item The current results (analytic free energy, geometric interpretation, numerical landscapes) are fully sufficient for a workshop submission.
    \item Extension toward ICML main track would require:
    \begin{itemize}
        \item clear identification of phase transitions,
        \item comparison across kernels,
        \item explicit characterization of retrieval failure regimes.
    \end{itemize}
\end{itemize}

%==============================================================================
\section*{Key Takeaways}
%==============================================================================

\begin{enumerate}
    \item The work studies thermodynamic limits of associative memory, not retrieval algorithms.
    \item Temperature represents intrinsic instability of the memory substrate.
    \item Retrieval failure corresponds to a geometric phase transition in the free-energy landscape.
    \item Modern high-capacity DAMs necessitate local geometric analysis.
    \item The direction is scientifically sound and clearly distinct from existing approaches.
\end{enumerate}

\end{document}
